\subsection{فعالیت ایستادن}

\subsubsection{ماهیت فعالیت و منطق تحلیل}

فعالیت ایستادن از نظر ظاهری یک فعالیت ساده و کم‌دامنه است، اما از دید تحلیل سیگنال، یکی از آزمون‌های مهم برای بررسی کنترل وضعیت بدن محسوب می‌شود. در این فعالیت، اختلاف بین سطوح عملکردی معمولاً به صورت حرکت‌های بزرگ و آشکار ظاهر نمی‌شود، بلکه بیشتر در میزان نوسان بدن، پایداری سر، تغییرات اندام‌های تحتانی و الگوهای جبرانی کوچک دیده می‌شود. بنابراین، انتظار نمی‌رود فقط یک ویژگی منفرد بتواند کلاس‌ها را از یکدیگر جدا کند؛ بلکه ترکیبی از ویژگی‌های آماری استخراج‌شده از چند حسگر باید الگوی کلی پایداری و کنترل وضعیت را توصیف کند.

در این فعالیت، تحلیل نتایج با سه پرسش اصلی انجام شد. نخست این‌که آیا سیگنال‌های ثبت‌شده در وضعیت ایستاده دارای تغییرات قابل استفاده برای طبقه‌بندی هستند یا خیر. دوم این‌که آیا فضای ویژگی استخراج‌شده، با وجود تعداد زیاد ویژگی‌ها، دارای افزونگی است و می‌تواند با روش
\lr{PCA}
کاهش یابد یا خیر. سوم این‌که کدام مدل یادگیری ماشین می‌تواند با پایداری بیشتری سطوح عملکردی را تفکیک کند.

\subsubsection{رفتار سیگنال و توزیع داده‌ها}

نمونه سیگنال فعالیت ایستادن در
\autoref{fig:stand_normal_signal_sample}
آورده شده است. برخلاف فعالیت‌های دینامیک، در این آزمون تغییرات سیگنال نباید به صورت ضربه‌ای یا جهشی باشد؛ بلکه دامنه‌های کوچک و نوسان‌های پیوسته اهمیت بیشتری دارند. همین موضوع باعث می‌شود که استخراج ویژگی‌های آماری مانند میانگین، انحراف معیار، دامنه تغییرات، انرژی، شاخص‌های مشتق‌شده و ویژگی‌های وابسته به تغییرات زمانی برای این فعالیت ضروری باشد. به عبارت دیگر، در ایستادن، اطلاعات تشخیصی بیشتر در «الگوی پایداری» قرار دارد، نه در یک رویداد حرکتی مشخص.

\begin{figure}[H]
	\centering
	\includegraphics[width=0.90\linewidth]{Images/Generated/TaskResults/Stand/Normal/signal_sample.png}
	\caption{نمونه سیگنال‌های ثبت‌شده در فعالیت ایستادن}
	\label{fig:stand_normal_signal_sample}
\end{figure}

توزیع کلاس‌ها در
\autoref{fig:stand_normal_class_distribution}
برای تفسیر نتایج مدل اهمیت دارد. اگر تعداد نمونه‌های سه کلاس تقریباً متعادل باشد، معیارهای کلی مانند
\lr{Accuracy}
قابل اعتمادتر تفسیر می‌شوند. با این حال، حتی در صورت تعادل داده‌ها نیز معیار
\lr{Accuracy}
برای این مسئله کافی نیست؛ زیرا ممکن است مدل در تشخیص یک کلاس خاص ضعیف عمل کند، ولی به دلیل عملکرد بهتر در کلاس‌های دیگر، میانگین کلی قابل قبول به نظر برسد. به همین دلیل، در ادامه ماتریس درهم‌ریختگی و
\lr{Recall}
هر کلاس نیز بررسی شده‌اند.

\begin{figure}[H]
	\centering
	\includegraphics[width=0.70\linewidth]{Images/Generated/TaskResults/Stand/Normal/class_distribution.png}
	\caption{توزیع نمونه‌ها در کلاس‌های مختلف فعالیت ایستادن}
	\label{fig:stand_normal_class_distribution}
\end{figure}

\subsubsection{افزونگی ویژگی‌ها و ضرورت کاهش بعد}

در این پروژه برای هر نمونه تعداد زیادی ویژگی از حسگرها و محورهای مختلف استخراج شده است. تعداد زیاد ویژگی‌ها از یک طرف باعث می‌شود اطلاعات متنوعی از سیگنال‌ها در اختیار مدل قرار گیرد، اما از طرف دیگر می‌تواند باعث افزونگی، همبستگی بالا و افزایش پیچیدگی مدل شود. نقشه همبستگی
\autoref{fig:stand_normal_feature_correlation}
نشان می‌دهد که بخشی از ویژگی‌ها رفتار مشابهی دارند و در خوشه‌های همبسته قرار می‌گیرند. این موضوع از نظر تحلیلی مهم است، زیرا نشان می‌دهد استفاده از تمام ویژگی‌های اولیه الزاماً به معنی استفاده از اطلاعات کاملاً مستقل نیست.

بنابراین، تحلیل
\lr{PCA}
در این فعالیت فقط یک مرحله جانبی نیست، بلکه روشی برای بررسی این پرسش است که چه مقدار از اطلاعات مفید داده را می‌توان با تعداد کمتری مؤلفه حفظ کرد. اگر تعداد مؤلفه‌های مورد نیاز برای حفظ ۹۵ درصد واریانس به طور قابل توجهی کمتر از تعداد ویژگی‌های اولیه باشد، می‌توان نتیجه گرفت که فضای ویژگی دارای افزونگی قابل توجه است.

% TODO: مقدار دقیق تعداد مؤلفه لازم برای رسیدن به ۹۵ درصد واریانس را از نمودار یا فایل Excel وارد کن.

\begin{figure}[H]
	\centering
	\includegraphics[width=0.90\linewidth]{Images/Generated/TaskResults/Stand/Normal/feature_correlation_heatmap.png}
	\caption{نقشه همبستگی ویژگی‌های منتخب در فعالیت ایستادن}
	\label{fig:stand_normal_feature_correlation}
\end{figure}

\autoref{fig:stand_normal_pca_variance}
نشان می‌دهد که واریانس داده چگونه بین مؤلفه‌های اصلی توزیع شده است. شیب اولیه نمودار برای تصمیم‌گیری اهمیت دارد: اگر چند مؤلفه اول بخش بزرگی از واریانس را توضیح دهند، کاهش بعد می‌تواند بدون از دست دادن بخش عمده اطلاعات انجام شود. در مقابل، اگر واریانس به صورت یکنواخت بین مؤلفه‌های زیاد پخش شده باشد، کاهش شدید بعد ممکن است بخشی از اطلاعات مؤثر در طبقه‌بندی را حذف کند.

\begin{figure}[H]
	\centering
	\includegraphics[width=0.90\linewidth]{Images/Generated/TaskResults/Stand/Normal/pca_variance_curve.png}
	\caption{واریانس توضیح‌داده‌شده و واریانس تجمعی مؤلفه‌های اصلی در فعالیت ایستادن}
	\label{fig:stand_normal_pca_variance}
\end{figure}

پراکندگی دوبعدی داده‌ها در فضای دو مؤلفه اول در
\autoref{fig:stand_normal_pca_scatter}
آورده شده است. این نمودار برای قضاوت نهایی کافی نیست، زیرا فقط دو مؤلفه اول را نمایش می‌دهد، اما برای شناخت اولیه ساختار داده مفید است. اگر کلاس‌ها در این فضا هم‌پوشانی داشته باشند، به این معنا نیست که طبقه‌بندی ناممکن است؛ بلکه نشان می‌دهد مدل برای تصمیم‌گیری به مؤلفه‌های بعدی یا ترکیب غیرخطی ویژگی‌ها نیاز دارد. این موضوع به‌ویژه برای فعالیت ایستادن قابل انتظار است، زیرا تفاوت بین سطح‌های عملکردی در این فعالیت ظریف‌تر از فعالیت‌های دینامیک است.

\begin{figure}[H]
	\centering
	\includegraphics[width=0.72\linewidth]{Images/Generated/TaskResults/Stand/Normal/pca_scatter_pc1_pc2.png}
	\caption{پراکندگی نمونه‌های فعالیت ایستادن در فضای دو مؤلفه اصلی اول}
	\label{fig:stand_normal_pca_scatter}
\end{figure}

\subsubsection{رفتار مدل‌ها و اثر کاهش بعد}

پس از بررسی ساختار ویژگی‌ها، سه مدل
\lr{KNN}،
\lr{SVM}
و
\lr{Random Forest}
روی داده‌های فعالیت ایستادن ارزیابی شدند. مقایسه کلی عملکرد مدل‌ها در
\autoref{fig:stand_normal_model_comparison_f1}
ارائه شده است. این نمودار باید همراه با نمودار اثر
\lr{PCA}
تفسیر شود، زیرا هدف فقط پیدا کردن بیشترین مقدار
\lr{F1}
نیست؛ بلکه باید مشخص شود آیا کاهش بعد می‌تواند با حفظ عملکرد، پیچیدگی فضای ویژگی را کمتر کند یا خیر.

اگر در این فعالیت مقدار
\lr{F1}
پس از اعمال
\lr{PCA95}
افزایش یافته باشد، می‌توان نتیجه گرفت که حذف افزونگی ویژگی‌ها به مدل کمک کرده است. اگر عملکرد تقریباً ثابت مانده باشد، نتیجه مهم‌تر این است که با تعداد بسیار کمتری مؤلفه می‌توان عملکرد مشابهی به دست آورد. اگر افت عملکرد مشاهده شود، باید بررسی شود که آیا افت ایجادشده در برابر کاهش تعداد ویژگی‌ها قابل قبول است یا خیر.

% TODO: در این پاراگراف مقدار دقیق تغییر F1 قبل و بعد از PCA را از شکل یا Excel اضافه کن.

\begin{figure}[H]
	\centering
	\includegraphics[width=0.88\linewidth]{Images/Generated/TaskResults/Stand/Normal/model_comparison_f1.png}
	\caption{مقایسه عملکرد مدل‌ها بر اساس معیار \lr{F1} در فعالیت ایستادن}
	\label{fig:stand_normal_model_comparison_f1}
\end{figure}

رتبه‌بندی مدل‌ها در
\autoref{fig:stand_normal_model_ranking_f1}
نشان می‌دهد کدام مدل در حالت انتخاب‌شده عملکرد بهتری داشته است. با این حال، مدل منتخب نهایی نباید فقط بر اساس این نمودار تعیین شود. برای یک سامانه قابل استفاده، پایداری مدل در تقسیم‌های مختلف داده و رفتار آن در هر کلاس نیز اهمیت دارد.

\begin{figure}[H]
	\centering
	\includegraphics[width=0.72\linewidth]{Images/Generated/TaskResults/Stand/Normal/model_ranking_f1.png}
	\caption{رتبه‌بندی مدل‌ها در فعالیت ایستادن بر اساس معیار \lr{F1}}
	\label{fig:stand_normal_model_ranking_f1}
\end{figure}

نمودار
\autoref{fig:stand_normal_f1_gain_pca}
به طور مستقیم نشان می‌دهد که کاهش بعد برای هر مدل چه اثری داشته است. این نمودار از نظر تصمیم‌گیری مهم‌تر از مقایسه خام مدل‌ها است، زیرا مشخص می‌کند آیا
\lr{PCA}
برای هر الگوریتم نقش مثبت، خنثی یا منفی داشته است. برای مثال، ممکن است یک مدل پس از
\lr{PCA}
کمی افت کند، اما به دلیل کاهش زیاد تعداد ویژگی‌ها همچنان انتخاب مناسبی باشد. در مقابل، اگر افت عملکرد زیاد باشد، استفاده از ویژگی‌های اولیه یا انتخاب تعداد مؤلفه متفاوت باید بررسی شود.

\begin{figure}[H]
	\centering
	\includegraphics[width=0.72\linewidth]{Images/Generated/TaskResults/Stand/Normal/f1_gain_pca_vs_no_pca.png}
	\caption{تغییر معیار \lr{F1} پس از اعمال \lr{PCA95} در فعالیت ایستادن}
	\label{fig:stand_normal_f1_gain_pca}
\end{figure}

پراکندگی معیارها در
\autoref{fig:stand_normal_metric_boxplots}
برای بررسی پایداری مدل‌ها استفاده شد. اگر یک مدل میانگین عملکرد بالایی داشته باشد اما پراکندگی زیادی بین
\lr{fold}
ها نشان دهد، نتیجه آن نسبت به تقسیم داده حساس است. در مقابل، مدلی که میانگین مناسب و پراکندگی کمتر داشته باشد، برای انتخاب نهایی قابل اتکاتر است.

\begin{figure}[H]
	\centering
	\includegraphics[width=0.90\linewidth]{Images/Generated/TaskResults/Stand/Normal/metric_boxplots.png}
	\caption{پراکندگی معیارهای ارزیابی در \lr{fold}های مختلف برای فعالیت ایستادن}
	\label{fig:stand_normal_metric_boxplots}
\end{figure}

% TODO: عددهای این جدول را از فایل‌های Excel مربوط به Stand/Normal وارد کن.
\begin{table}[H]
	\centering
	\caption{خلاصه عملکرد مدل‌ها برای فعالیت ایستادن}
	\label{tab:stand_model_metrics}
	\begin{tabular}{lccccc}
		\toprule
		مدل & \lr{Accuracy} & \lr{Precision} & \lr{Recall} & \lr{Specificity} & \lr{F1} \\
		\midrule
		\lr{KNN} &  &  &  &  &  \\
		\lr{SVM} &  &  &  &  &  \\
		\lr{Random Forest} &  &  &  &  &  \\
		\bottomrule
	\end{tabular}
\end{table}

\subsubsection{تحلیل خطاهای طبقه‌بندی}

برای بررسی دقیق‌تر رفتار مدل‌ها، ماتریس‌های درهم‌ریختگی نرمال‌شده در دو حالت بدون
\lr{PCA}
و با
\lr{PCA95}
در
\autoref{fig:stand_normal_confusion_no_pca}
و
\autoref{fig:stand_normal_confusion_pca95}
ارائه شده‌اند. در فعالیت ایستادن، خطا بین کلاس‌های مجاور از نظر حرکتی قابل انتظارتر است، زیرا تفاوت بین اجرای اصلاح‌شده و اجرای طبیعی ممکن است بسیار ظریف باشد. بنابراین اگر بخش عمده خطاها بین کلاس‌های نزدیک رخ داده باشد، این موضوع بیشتر به دشواری ذاتی فعالیت مربوط است تا ضعف کامل مدل.

\begin{figure}[H]
	\centering
	\includegraphics[width=0.92\linewidth]{Images/Generated/TaskResults/Stand/Normal/confusion_matrices_no_pca.png}
	\caption{ماتریس‌های درهم‌ریختگی نرمال‌شده در فعالیت ایستادن، بدون \lr{PCA}}
	\label{fig:stand_normal_confusion_no_pca}
\end{figure}

\begin{figure}[H]
	\centering
	\includegraphics[width=0.92\linewidth]{Images/Generated/TaskResults/Stand/Normal/confusion_matrices_pca95.png}
	\caption{ماتریس‌های درهم‌ریختگی نرمال‌شده در فعالیت ایستادن، با \lr{PCA95}}
	\label{fig:stand_normal_confusion_pca95}
\end{figure}

نمودارهای
\lr{Recall}
هر کلاس در
\autoref{fig:stand_normal_per_class_recall_no_pca}
و
\autoref{fig:stand_normal_per_class_recall_pca95}
نشان می‌دهند که آیا مدل برای یک سطح عملکردی خاص ضعف دارد یا خیر. این بخش برای کاربرد بالینی اهمیت زیادی دارد، زیرا تشخیص اشتباه سطح عملکردی پایین یا متوسط می‌تواند تفسیر نهایی سامانه را تغییر دهد. بنابراین، حتی اگر میانگین
\lr{F1}
مناسب باشد، باید اطمینان حاصل شود که هیچ کلاس مهمی به طور سیستماتیک نادیده گرفته نشده است.

\begin{figure}[H]
	\centering
	\includegraphics[width=0.86\linewidth]{Images/Generated/TaskResults/Stand/Normal/per_class_recall_no_pca.png}
	\caption{مقایسه \lr{Recall} هر کلاس در فعالیت ایستادن، بدون \lr{PCA}}
	\label{fig:stand_normal_per_class_recall_no_pca}
\end{figure}

\begin{figure}[H]
	\centering
	\includegraphics[width=0.86\linewidth]{Images/Generated/TaskResults/Stand/Normal/per_class_recall_pca95.png}
	\caption{مقایسه \lr{Recall} هر کلاس در فعالیت ایستادن، با \lr{PCA95}}
	\label{fig:stand_normal_per_class_recall_pca95}
\end{figure}

\subsubsection{اهمیت ویژگی‌ها و انتخاب تعداد مؤلفه‌ها}

تحلیل اهمیت ویژگی‌ها در
\autoref{fig:stand_normal_rf_importance}
به فهم این موضوع کمک می‌کند که مدل بیشتر به کدام حسگرها، محورهای حرکتی یا شاخص‌های آماری تکیه کرده است. اگر ویژگی‌های مرتبط با سر یا پاها سهم بیشتری داشته باشند، این نتیجه با ماهیت فعالیت ایستادن سازگار است؛ زیرا پایداری سر و کنترل اندام‌های تحتانی در حفظ وضعیت ایستاده نقش دارند. اگر ویژگی‌های دست نیز در میان ویژگی‌های مهم ظاهر شوند، می‌توان آن را به حرکات جبرانی یا نوسانات همزمان اندام فوقانی نسبت داد.

\begin{figure}[H]
	\centering
	\includegraphics[width=0.90\linewidth]{Images/Generated/TaskResults/Stand/Normal/rf_top_feature_importance.png}
	\caption{ویژگی‌های مهم‌تر بر اساس مدل \lr{Random Forest} در فعالیت ایستادن}
	\label{fig:stand_normal_rf_importance}
\end{figure}

نقشه حرارتی میانگین نرمال‌شده ویژگی‌ها در
\autoref{fig:stand_normal_class_centroid_heatmap}
برای بررسی تفاوت الگوی ویژگی‌ها بین کلاس‌ها استفاده شد. اگر کلاس‌ها الگوی رنگی متفاوتی داشته باشند، نشان می‌دهد که تفاوت بین سطح‌های عملکردی فقط در خروجی مدل ظاهر نشده، بلکه در خود فضای ویژگی نیز قابل مشاهده است.

\begin{figure}[H]
	\centering
	\includegraphics[width=0.90\linewidth]{Images/Generated/TaskResults/Stand/Normal/class_centroid_feature_heatmap.png}
	\caption{میانگین نرمال‌شده ویژگی‌های منتخب در کلاس‌های مختلف فعالیت ایستادن}
	\label{fig:stand_normal_class_centroid_heatmap}
\end{figure}

در نهایت، اثر تعداد مؤلفه‌های اصلی بر عملکرد مدل‌ها در
\autoref{fig:stand_normal_svm_metrics_vs_pca_components}
و
\autoref{fig:stand_normal_models_f1_vs_pca_components}
بررسی شده است. این بخش برای انتخاب تنظیمات نهایی اهمیت دارد. اگر عملکرد مدل از یک تعداد مؤلفه به بعد تقریباً ثابت شود، استفاده از مؤلفه‌های بیشتر فقط پیچیدگی را افزایش می‌دهد. اگر با کاهش مؤلفه‌ها افت شدید رخ دهد، باید تعداد مؤلفه بیشتری نگه داشته شود یا انتخاب ویژگی با روش دیگری انجام شود.

% TODO: پس از خواندن نمودار، تعداد مؤلفه‌ای را که بهترین F1 یا بهترین مصالحه بین F1 و کاهش بعد ایجاد کرده است، در متن وارد کن.

\begin{figure}[H]
	\centering
	\includegraphics[width=0.88\linewidth]{Images/Generated/TaskResults/Stand/Normal/svm_metrics_vs_pca_components.png}
	\caption{تغییر عملکرد مدل \lr{SVM} نسبت به تعداد مؤلفه‌های اصلی در فعالیت ایستادن}
	\label{fig:stand_normal_svm_metrics_vs_pca_components}
\end{figure}

\begin{figure}[H]
	\centering
	\includegraphics[width=0.88\linewidth]{Images/Generated/TaskResults/Stand/Normal/models_f1_vs_pca_components.png}
	\caption{مقایسه مدل‌ها نسبت به تعداد مؤلفه‌های اصلی در فعالیت ایستادن}
	\label{fig:stand_normal_models_f1_vs_pca_components}
\end{figure}

% TODO: پس از انتخاب نهایی مدل، مقدار مؤلفه و معیارهای نهایی این جدول را پر کن.
\begin{table}[H]
	\centering
	\caption{تنظیمات منتخب فعالیت ایستادن برای مدل نهایی}
	\label{tab:stand_selected_setup}
	\begin{tabular}{lccc}
		\toprule
		مدل منتخب & نوع داده & تعداد مؤلفه اصلی & معیار انتخاب \\
		\midrule
		 & داده‌های اصلی &  &  \\
		\bottomrule
	\end{tabular}
\end{table}

\subsubsection{جمع‌بندی فعالیت ایستادن}

جمع‌بندی فعالیت ایستادن نشان می‌دهد که با وجود کم‌دامنه بودن حرکت، داده‌های حسگرهای اینرسیایی می‌توانند اطلاعات قابل استفاده‌ای برای تفکیک سطوح عملکردی فراهم کنند. نکته اصلی در این فعالیت، تشخیص الگوهای ظریف پایداری و نوسان است. بنابراین، انتخاب مدل نهایی باید فقط بر اساس بیشترین مقدار میانگین انجام نشود، بلکه لازم است اثر کاهش بعد، پایداری در
\lr{fold}
ها، خطاهای بین‌کلاسی و اهمیت ویژگی‌ها به صورت همزمان در نظر گرفته شود.
